% =============================================================================
% KAPITEL 7: DISKUSSION DER ERGEBNISSE (INKL. LIMITATIONEN)
% =============================================================================

\chapter{Diskussion der Ergebnisse}\label{chap:diskussion}

Dieses Kapitel ordnet die in Kapitel~\ref{chap:evaluation} präsentierten
Befunde in den wissenschaftlichen Kontext ein, diskutiert ihre Tragweite
und benennt die Grenzen der gewählten Methodik. Abschnitt~\ref{sec:disk-befunde}
diskutiert die zentralen inhaltlichen Befunde. Abschnitt~\ref{sec:disk-limitationen}
benennt die methodischen Limitationen der Arbeit. Abschnitt~\ref{sec:disk-zusammenfassung}
schließt mit einer Gesamtbewertung.

% =============================================================================
\section{Diskussion der zentralen Befunde}
\label{sec:disk-befunde}
% =============================================================================

\subsection{xG als prädiktiver Mehrwert -- Konsistenz und Grenzen}

Der wichtigste Befund dieser Arbeit ist, dass die Integration von
Expected-Goals-Daten den F1-macro-Score in allen vier paarweisen Vergleichen
verbessert. Der mittlere Gewinn von $+0{,}041$ F1-Punkten ist moderat,
aber konsistent -- und genau diese Konsistenz ist wissenschaftlich
belastbarer als ein einzelner hoher Ausreißer in einer Konfiguration.
Die Ergebnisse stehen im Einklang mit \textcite{AnzerBauer2021} und
\textcite{VanEetveldeEtAl2024}, die xG als stabilen Teamleistungsindikator
identifiziert haben.

Einschränkend ist jedoch festzuhalten, dass die verwendeten xG-Werte
\emph{saisonale Durchschnitte pro Spiel} sind -- also ein einziger
statischer Wert für alle Spiele eines Teams in einer gesamten Saison.
\textcite{ForcherEtAl2025} haben für die Bundesliga gezeigt, dass
matchweises Post-Match-xG eine Accuracy von 0,656 ermöglicht. Der in
dieser Arbeit erzielte Accuracy-Wert von 0,499 liegt deutlich darunter
-- was primär auf den gröberen Zeitmaßstab der Pre-Match-xG-Daten
zurückzuführen ist. Der prädiktive Mehrwert von xG wäre vermutlich
größer, wenn matchweise oder spielphasenweise xG-Daten verfügbar wären
\cite{Wheatcroft2022}. Das Modell nutzt damit nicht das volle Potential
des xG-Konzepts, sondern nur einen saisonalen Proxy -- der dennoch
einen messbaren Informationsgewinn liefert.

\subsection{Random Forest als bestes Modell -- eine überraschende Erkenntnis}

Entgegen der in der Literatur häufig beschriebenen Überlegenheit von
Gradient-Boosting-Verfahren bei tabellarischen Daten \cite{ChenGuestrin2016,KeEtAl2017}
erzielt in dieser Arbeit der Random Forest (F1-macro 0,479) den besten
Wert, gefolgt von LightGBM und XGBoost (je 0,473). Dieser Befund ist
plausibel erklärbar: Gradient Boosting profitiert typischerweise von
größeren Datensätzen, bei denen die sequenzielle Residuenkorrektur über
viele Iterationen hinweg ihre volle Wirkung entfalten kann. Mit
$\sim$2.300 Trainingsbeispielen in der besten xG-Variante ist die
Datenmenge begrenzt, was das Bagging-Prinzip des Random Forest -- das
Rauschen durch Mittelung vieler unabhängig trainierter Bäume reduziert
-- begünstigt. Dieses Ergebnis reiht sich in Beobachtungen ein, die
zeigen, dass auf kleinen Datensätzen Ensemble-Methoden unterschiedlicher
Art vergleichbar performen und keine pauschale Überlegenheit eines
einzelnen Verfahrens existiert \cite{BunkerThabtah2019}.

\subsection{Unentschieden-Erkennung -- Fortschritt und verbleibende Grenzen}

Alle fünf Modelle erreichen F1-Draw-Werte zwischen 0,305 und 0,385 und
übertreffen damit die in der Literatur für dieses Problem typische
Bandbreite von 0,10--0,30 \cite{BunkerThabtah2019,HubacekSourekZelezny2019}.
Dies zeigt, dass die eingesetzten Gegenmaßnahmen -- klassengewichtetes
Training und manuelles Oversampling beim MLP -- ihre beabsichtigte
Wirkung entfaltet haben. Der gezielt konstruierte
\texttt{combined\_draw\_tendency}-Index trägt ebenfalls zur
Unentschieden-Erkennung bei, ohne jedoch als einzelnes dominantes
Feature hervorzustechen.

Dennoch bleibt der F1-Draw mit 0,331 beim Random Forest deutlich hinter
dem F1-Heimsieg zurück. Dies ist kein Versagen der Modellierung, sondern
spiegelt eine fundamentale epistemische Grenze wider: Unentschieden
entstehen in Spielen, in denen beide Teams annähernd gleich stark sind
-- eine Ausgangslage, die mit pre-match aggregierten Formdaten nur
unzureichend zu quantifizieren ist, da in solchen Partien kleinste
zufällige Ereignisse den Ausschlag geben \cite{DobsonGoddard2011,Sumpter2016}.

% =============================================================================
\section{Limitationen der Arbeit}
\label{sec:disk-limitationen}
% =============================================================================

\subsection{Fehlende Einflussvariablen}

Selbst das beste entwickelte Modell erfasst nur einen Bruchteil der
spielrelevanten Informationen. Folgende Faktoren sind in den Daten
nicht enthalten, beeinflussen Spielergebnisse aber nachweislich:

\begin{itemize}
    \item \textbf{Verletzungen und Kaderabsenzen:} Der Ausfall von
          Schlüsselspielern kann den Erwartungswert eines Spiels erheblich
          verschieben \cite{DobsonGoddard2011}. EWA-Features mitteln
          über Perioden mit und ohne diese Informationen hinweg.
    \item \textbf{Rote Karten und Spielunterbrechungen:} Solche Ereignisse
          verändern den Spielverlauf fundamental und können nicht antizipiert
          werden, da sie stochastisch auftreten.
    \item \textbf{Trainerwechsel und taktische Umstellungen:} Nach einem
          Trainerwechsel können EWA-Formwerte systematisch falsche Signale
          senden, da sie auf dem früheren Spielstil basieren.
    \item \textbf{Motivation und Kontextfaktoren:} Ein Team, das bereits
          den Abstieg besiegelt hat, verhält sich anders als ein Team im
          Titelkampf -- eine Information, die keiner der verwendeten
          Features kodiert.
\end{itemize}

Diese fehlenden Variablen sind Teil des in Kapitel~\ref{chap:grundlagen}
beschriebenen irreduziblen Fehlers $\sigma^2$: Sie erhöhen die prinzipielle
Unvorhersehbarkeit von Fußballergebnissen und begrenzen die erreichbare
Modellgüte unabhängig vom eingesetzten Algorithmus \cite{GemanEtAl1992}.

\subsection{Granularität und Aktualität der xG-Daten}

Die von \textcite{Understat2024} bezogenen xG-Daten sind saisonale
Aggregatwerte: Für jedes Team gibt es pro Saison genau einen xG-Wert
pro Spiel, der für alle 34 Saisonspiele identisch ist. Damit fehlt
die zeitliche Dynamik innerhalb einer Saison vollständig. Das Lag-1-Design
-- Vorsaison-xG als Prior für die aktuelle Saison -- ist methodisch
korrekt und verhindert Data Leakage \cite{KaufmanRossetPerlich2012},
bedeutet aber auch, dass die xG-Werte immer mindestens eine Saison
\emph{veraltet} sind, wenn sie in die Vorhersage einfließen. Ein Team
wie Bayer Leverkusen, das nach einem starken xG-Jahr in der Folgesaison
die Kaderqualität massiv verändert, wird durch den Lag-1-Prior falsch
eingeschätzt.

\subsection{Zeitliche Validität des Datensplits}

Der stratifizierte 80:20-Split teilt den Datensatz zufällig auf, ohne
die zeitliche Struktur der Daten zu berücksichtigen. In der Praxis
würde ein Vorhersagemodell ausschließlich auf vergangenen Spielen
trainiert und auf zukünftigen Spielen getestet. Der hier gewählte
zufällige Split erlaubt es theoretisch, dass Trainingsdaten aus der
Saison 2024/25 zur Vorhersage von Spielen aus der Saison 2018/19
genutzt werden -- eine zeitliche Inkonsistenz, die in realen
Anwendungsszenarien nicht möglich wäre. Eine temporal-konsistente
Aufteilung (z.\,B.\ Training auf 2016/17--2022/23, Test auf 2023/24
bis 2025/26) würde die Realitätsnähe des Evaluationsdesigns verbessern,
birgt jedoch das Risiko kleinerer und weniger repräsentativer Testsets.

\subsection{Eingeschränkte Hyperparameter-Optimierung}

Die Randomized Search mit 10--20 Iterationen je Modell ist eine
pragmatische Annäherung an die Hyperparameter-Optimierung. Bei einem
Suchraum von mehreren tausend Kombinationen pro Algorithmus besteht
die Möglichkeit, dass deutlich bessere Konfigurationen in den
durchgeführten Iterationen nicht gezogen wurden. Die beobachtete
CV-Stabilität des Random-Forest-Modells mit einer Standardabweichung
von $\pm$0,014 ist allerdings ein Indiz dafür, dass die Performance
nicht stark hyperparameter-sensitiv ist -- die Grundaussage der Ergebnisse
dürfte daher robust sein \cite{JamesEtAl2021}.

\subsection{Fehlende Bezeichnung der Feature-Importance}

Da die Feature-Importance-Analyse im vorliegenden Code generische
Bezeichnungen (F0--F29) statt der echten Feature-Namen verwendet, ist
die Zuordnung der wichtigsten Features zu konkreten Metriken nur durch
Indexrückschluss möglich. Eine präzisere Implementierung würde die
Featurenamen direkt aus den verwendeten Listen übergeben, was die
Interpretierbarkeit der Modelle erheblich verbessern würde
\cite{LundbergLee2017}.

% =============================================================================
\section{Gesamtbewertung}
\label{sec:disk-zusammenfassung}
% =============================================================================

Die vorliegende Arbeit zeigt, dass Machine-Learning-basierte Vorhersage
von Bundesliga-Spielergebnissen mit sorgfältigem Feature Engineering
und gezielter Behandlung der Klassen-Imbalance zu Ergebnissen führt,
die den in der Literatur beschriebenen Stand des Möglichen für Pre-Match-Vorhersagen
widerspiegeln. Die Integration von Expected-Goals-Daten liefert einen
konsistenten, wenn auch moderaten Mehrwert.

Die methodischen Stärken der Arbeit liegen in der konsequenten
Vermeidung von Data Leakage durch den Lag-1-Ansatz und die Shift-Operation
bei der EWA-Berechnung, in der systematischen Evaluation von zwölf
Datensatzvarianten über fünf Algorithmen, sowie in der gezielten
Konstruktion von Unentschieden-Features, die nachweislich zu überdurchschnittlichen
F1-Draw-Werten führen.

Die wesentliche Limitation liegt in der Granularität der verfügbaren
xG-Daten: Saisonale Pre-Match-Durchschnitte liefern ein deutlich
schwächeres Signal als matchweises Post-Match-xG. Zukünftige Arbeiten
sollten matchweise xG-Zeitreihen verwenden und ein temporal-konsistentes
Evaluationsdesign wählen, um die Übertragbarkeit der Befunde auf
reale Vorhersageszenarien zu stärken.